\section{Conclusion}\label{sec:conclusion}

This paper formulates a proof architecture for quantized GBDT in which the
central training difficulty, namely optimal split selection, is handled through
a dedicated primitive layer rather than through a monolithic generic circuit.
By committing the training trace as structured tensors, \textsc{Terrae} proves
routing, histogram construction, split optimality, leaf-value assignment, and
boosting-state updates directly over polynomial commitments. Its main technical
tools, domain-lifting batching and interleaving batching, reduce many
heterogeneous algebraic and non-native checks to a small number of batched
claims without introducing extra auxiliary commitments. The histogram
aggregation argument further avoids a generic RAM proof for the dominant
sample-to-bucket update pattern in tree training. Together, these components
give a specialized zero-knowledge proof system for GBDT training and inference,
and the evaluation shows that exploiting this structure substantially reduces
prover cost relative to a circuit-style baseline.
